% ═══════════════════════════════════════════════════════════════
% LEHRERHINWEISE: Pretérito Indefinido (Einführung -ar Verben)
% ═══════════════════════════════════════════════════════════════
% Sequenz: 08 - Las vacaciones
% Stunde: 08c - Pretérito Indefinido Einführung
% ═══════════════════════════════════════════════════════════════

\documentclass[11pt, a4paper]{article}

% ═══════════════════════════════════════════════════════════════
% SW-MODUS TOGGLE
% ═══════════════════════════════════════════════════════════════
\newif\ifswmodus
\swmodusfalse

% ═══════════════════════════════════════════════════════════════
% PAKETE
% ═══════════════════════════════════════════════════════════════

\usepackage[utf8]{inputenc}
\usepackage[T1]{fontenc}
\usepackage[ngerman]{babel}
\usepackage{helvet}
\renewcommand{\familydefault}{\sfdefault}

\usepackage[margin=2cm]{geometry}
\usepackage{enumitem}
\usepackage{tabularx}
\usepackage{booktabs}
\usepackage{longtable}

\usepackage{xcolor}
\usepackage{amssymb}
\usepackage{tcolorbox}
\tcbuselibrary{skins}

\usepackage{fancyhdr}
\usepackage{lastpage}
\usepackage{needspace}

% ═══════════════════════════════════════════════════════════════
% FARBEN
% ═══════════════════════════════════════════════════════════════

\definecolor{spanisch-color}{HTML}{c41e3a}
\definecolor{grau-color}{HTML}{888888}
\definecolor{hellgrau-color}{HTML}{f5f5f5}
\definecolor{basis-color}{HTML}{2a7a4b}
\definecolor{standard-color}{HTML}{2c5aa0}
\definecolor{erweiterung-color}{HTML}{7b1fa2}
\definecolor{loesung-color}{HTML}{c62828}

\definecolor{sw-dark}{HTML}{000000}
\definecolor{sw-gray}{HTML}{666666}
\definecolor{sw-white}{HTML}{ffffff}

\ifswmodus
    \colorlet{spanisch}{sw-dark}
    \colorlet{grau}{sw-gray}
    \colorlet{hellgrau}{sw-white}
    \colorlet{basis}{sw-dark}
    \colorlet{standard}{sw-dark}
    \colorlet{erweiterung}{sw-dark}
    \colorlet{loesung}{sw-dark}
\else
    \colorlet{spanisch}{spanisch-color}
    \colorlet{grau}{grau-color}
    \colorlet{hellgrau}{hellgrau-color}
    \colorlet{basis}{basis-color}
    \colorlet{standard}{standard-color}
    \colorlet{erweiterung}{erweiterung-color}
    \colorlet{loesung}{loesung-color}
\fi

\newcommand{\fachfarbe}{spanisch}

% ═══════════════════════════════════════════════════════════════
% VARIABLEN
% ═══════════════════════════════════════════════════════════════

\newcommand{\thema}{Pretérito Indefinido (Einführung -ar)}
\newcommand{\sequenz}{08}
\newcommand{\block}{c}
\newcommand{\fach}{Spanisch}
\newcommand{\klasse}{7}
\newcommand{\dauer}{90 Min. (Doppelstunde)}
\newcommand{\lerngruppengroesse}{24}
\newcommand{\datum}{\today}

% ═══════════════════════════════════════════════════════════════
% KOPF- UND FUSSZEILE
% ═══════════════════════════════════════════════════════════════

\pagestyle{fancy}
\fancyhf{}
\renewcommand{\headrulewidth}{0.4pt}
\renewcommand{\footrulewidth}{0.4pt}

\lhead{\fach{} -- Klasse \klasse}
\chead{\textbf{LH-\sequenz\block{} (Lehrerhinweise)}}
\rhead{\datum}

\lfoot{}
\cfoot{}
\rfoot{Seite \thepage\ von \pageref{LastPage}}

% ═══════════════════════════════════════════════════════════════
% BEFEHLE
% ═══════════════════════════════════════════════════════════════

\newcommand{\B}{\textcolor{basis}{\textbf{[B]}}}
\newcommand{\Std}{\textcolor{standard}{\textbf{[S]}}}
\newcommand{\E}{\textcolor{erweiterung}{\textbf{[E]}}}

\newcommand{\loesung}[1]{\textcolor{loesung}{\textbf{#1}}}

\newcommand{\es}[1]{\textcolor{\fachfarbe}{\textbf{#1}}}

\newtcolorbox{kurzplan}{
    colback=hellgrau,
    colframe=\fachfarbe,
    colbacktitle=white,
    coltitle=\fachfarbe,
    boxrule=1pt,
    arc=0pt,
    fonttitle=\bfseries,
    attach boxed title to top left={yshift=-2mm, xshift=5mm},
    boxed title style={boxrule=0pt, colback=white},
    title=Kurzplan
}

\newtcolorbox{differenzierung}{
    colback=white,
    colframe=grau,
    colbacktitle=white,
    coltitle=grau,
    boxrule=0.5pt,
    arc=0pt,
    fonttitle=\bfseries,
    attach boxed title to top left={yshift=-2mm, xshift=5mm},
    boxed title style={boxrule=0pt, colback=white},
    title=Differenzierung
}

\newtcolorbox{fehlerbox}{
    colback=white,
    colframe=loesung,
    colbacktitle=white,
    coltitle=loesung,
    boxrule=0.5pt,
    arc=0pt,
    fonttitle=\bfseries,
    attach boxed title to top left={yshift=-2mm, xshift=5mm},
    boxed title style={boxrule=0pt, colback=white},
    title=Häufige Fehler
}

\newtcolorbox{materialbox}{
    colback=white,
    colframe=\fachfarbe,
    colbacktitle=white,
    coltitle=\fachfarbe,
    boxrule=0.5pt,
    arc=0pt,
    fonttitle=\bfseries,
    attach boxed title to top left={yshift=-2mm, xshift=5mm},
    boxed title style={boxrule=0pt, colback=white},
    title=Material-Checkliste
}

% ═══════════════════════════════════════════════════════════════
% DOKUMENT
% ═══════════════════════════════════════════════════════════════

\begin{document}

% ═══════════════════════════════════════════════════════════════
% KOPFBEREICH
% ═══════════════════════════════════════════════════════════════

\begin{center}
    {\Large\bfseries\textcolor{\fachfarbe}{Lehrerhinweise: \thema}}\\[0.3em]
    {\normalsize LH-\sequenz\block{} | \dauer}
\end{center}

\vspace{10pt}

% ═══════════════════════════════════════════════════════════════
% 1. KURZPLAN
% ═══════════════════════════════════════════════════════════════

\section*{1. Stundenverlauf (Kurzplan)}

\textbf{1. Stunde (45 Min.)}

\begin{tabularx}{\textwidth}{|c|l|X|l|}
\hline
\textbf{Zeit} & \textbf{Phase} & \textbf{Inhalt} & \textbf{Material} \\
\hline
0--5 & Einstieg & Urlaubsfotos, Frage: ``¿Qué hicisteis en las vacaciones?'' & Beamer \\
\hline
5--12 & Input & Beispieltext, Verbformen markieren & Tafel, Text \\
\hline
12--22 & Erarbeitung & Induktive Regelfindung: Endungen entdecken & Tafel \\
\hline
22--30 & Systematisierung & Konjugationstabelle erstellen, Akzente erklären & Tafel, AB \\
\hline
30--38 & Übung 1 & Chorisches Sprechen: hablé, hablaste, habló... & -- \\
\hline
38--45 & Übung 2 & AB Aufgabe 1: Signalwörter zuordnen & AB \\
\hline
\end{tabularx}

\vspace{1em}

\textbf{2. Stunde (45 Min.)}

\begin{tabularx}{\textwidth}{|c|l|X|l|}
\hline
\textbf{Zeit} & \textbf{Phase} & \textbf{Inhalt} & \textbf{Material} \\
\hline
0--8 & Wiederholung & Konjugations-Quiz im Plenum & Tafel \\
\hline
8--20 & Übung 3 & AB Aufgaben 2--3: Formen bilden, Lückentext & AB \\
\hline
20--35 & Anwendung & AB Aufgabe 4: Eigener Urlaubsbericht schreiben & AB \\
\hline
35--42 & Sicherung & Partneraustausch, ausgewählte Texte vorlesen & -- \\
\hline
42--45 & Abschluss & Zusammenfassung, HA: Lernpfad + Vokabeln & Tafel \\
\hline
\end{tabularx}

% ═══════════════════════════════════════════════════════════════
% 2. DIFFERENZIERUNG
% ═══════════════════════════════════════════════════════════════

\section*{2. Differenzierung}

\begin{differenzierung}
\textbf{Legende:} \B{} Basis (BR) | \Std{} Standard (MR) | \E{} Erweiterung (GY)

\vspace{0.5em}

\textbf{WICHTIG:} Diese Marker sind NUR für die Lehrkraft. Auf Schülermaterial erscheinen KEINE Niveaumarkierungen!
\end{differenzierung}

\vspace{1em}

\begin{tabularx}{\textwidth}{|l|X|X|X|}
\hline
& \B{} \textbf{Basis} & \Std{} \textbf{Standard} & \E{} \textbf{Erweiterung} \\
\hline
\textbf{Aufgabe 1} & Mit Wortliste lösen & Selbstständig & + Sätze mit Signalwörtern \\
\hline
\textbf{Aufgabe 2} & Tabelle als Hilfe & Selbstständig & + Weitere Verben (komplett) \\
\hline
\textbf{Aufgabe 3} & Mit Endungsliste & Selbstständig & + Erweitern um eigene Sätze \\
\hline
\textbf{Aufgabe 4} & 3 Sätze mit Satzbaukasten & 4--5 Sätze & 6--8 Sätze, freier Text \\
\hline
\end{tabularx}

\vspace{1em}

\textbf{Konkrete Hinweise:}
\begin{itemize}
    \item \B{} SuS mit LRS: Konjugationstabelle vergrößert kopieren, Zeitzugabe +10 Min.
    \item \B{} Schwache SuS: Satzbaukasten als Hilfe bei Aufgabe 4
    \item \E{} Leistungsstarke: Expertenrolle beim Partneraustausch, Text auf Spanisch vorlesen
\end{itemize}

\textbf{Satzbaukasten für \B{}:}
\begin{center}
\fbox{\parbox{0.9\textwidth}{
\textbf{[Signalwort]} + \textbf{[Person]} + \textbf{[Verb im Indefinido]} + \textbf{[Rest des Satzes]}

\vspace{0.3em}

Beispiel: \es{El año pasado} + \es{yo} + \es{viajé} + \es{a España.}
}}
\end{center}

% ═══════════════════════════════════════════════════════════════
% 3. HÄUFIGE FEHLER
% ═══════════════════════════════════════════════════════════════

\section*{3. Häufige Fehler}

\begin{fehlerbox}
\begin{tabularx}{\textwidth}{|l|X|X|}
\hline
\textbf{Fehler} & \textbf{Beispiel} & \textbf{Reaktion} \\
\hline
Akzent vergessen & *hable statt \es{hablé} & Bedeutungsunterschied: hablo (Präsens) vs. habló (Vergangenheit). Visuelle Markierung an Tafel. \\
\hline
Falsche Endung & *hablí statt \es{hablé} & Auf Konjugationstabelle verweisen, chorisch wiederholen \\
\hline
Präsens statt Indefinido & *Ayer viajo a... & Signalwort ``ayer'' zeigt Vergangenheit! Immer Signalwort + Indefinido \\
\hline
-amos verwechselt & Präsens = Indefinido? & Ja, bei nosotros identisch! Kontext/Signalwort beachten \\
\hline
Infinitiv falsch erkannt & *nadé statt \es{nadé} & Stamm = Infinitiv ohne -ar. Dann Endung anhängen \\
\hline
\end{tabularx}
\end{fehlerbox}

\textbf{Typische Interferenz Deutsch:}
\begin{itemize}
    \item Deutsche SuS verwenden oft Perfekt (``Ich habe gespielt'') wo Spanisch Indefinido nutzt
    \item Hinweis: ``Ayer'' = abgeschlossen = Indefinido (nicht Perfecto)
\end{itemize}

% ═══════════════════════════════════════════════════════════════
% 4. MATERIAL-CHECKLISTE
% ═══════════════════════════════════════════════════════════════

\section*{4. Material-Checkliste}

\begin{materialbox}
\begin{itemize}[label=$\square$]
    \item AB-08c.pdf (\lerngruppengroesse{} Kopien)
    \item ML-08c.pdf (1x für Lehrkraft)
    \item Lehrwerk: Qué pasa 1, S. 140--144
    \item Urlaubsfotos/Bilder (Beamer oder Poster)
    \item Tafelmarker / Kreide
    \item Optional: LP-08c.html (Tablets/PCs für Hausaufgabe)
    \item Optional: Vergrößerte Konjugationstabelle für LRS-SuS
\end{itemize}
\end{materialbox}

% ═══════════════════════════════════════════════════════════════
% 5. LÖSUNGEN
% ═══════════════════════════════════════════════════════════════

\section*{5. Lösungen AB-08c}

\textbf{Merkkasten: Konjugation hablar}

\begin{tabular}{ll}
yo & \loesung{hablé} \\
tú & \loesung{hablaste} \\
él/ella/usted & \loesung{habló} \\
nosotros/as & \loesung{hablamos} \\
vosotros/as & \loesung{hablasteis} \\
ellos/ellas/ustedes & \loesung{hablaron} \\
\end{tabular}

\vspace{1em}

\textbf{Merkkasten: Signalwörter}

\begin{tabular}{ll}
\es{ayer} & $\rightarrow$ \loesung{gestern} \\
\es{el año pasado} & $\rightarrow$ \loesung{letztes Jahr} \\
\es{el verano pasado} & $\rightarrow$ \loesung{letzten Sommer} \\
\es{la semana pasada} & $\rightarrow$ \loesung{letzte Woche} \\
\end{tabular}

\vspace{1em}

\textbf{Aufgabe 1: Zuordnung} (4 Punkte)

\begin{tabular}{ll}
\es{ayer} & $\rightarrow$ \loesung{gestern} \\
\es{el año pasado} & $\rightarrow$ \loesung{letztes Jahr} \\
\es{la semana pasada} & $\rightarrow$ \loesung{letzte Woche} \\
\es{el verano pasado} & $\rightarrow$ \loesung{letzten Sommer} \\
\end{tabular}

\vspace{1em}

\textbf{Aufgabe 2: Verbformen} (6 Punkte)

\begin{tabular}{ll}
a) yo / hablar & \loesung{hablé} \\
b) tú / viajar & \loesung{viajaste} \\
c) él / nadar & \loesung{nadó} \\
d) nosotros / visitar & \loesung{visitamos} \\
e) vosotros / hablar & \loesung{hablasteis} \\
f) ellos / viajar & \loesung{viajaron} \\
\end{tabular}

\vspace{1em}

\textbf{Aufgabe 3: Lückentext} (6 Punkte)

\begin{tabular}{ll}
mi familia y yo (viajar) & \loesung{viajamos} \\
Yo (nadar) & \loesung{nadé} \\
Mi hermana (tomar) & \loesung{tomó} \\
Nosotros (visitar) & \loesung{visitamos} \\
Mis padres (hablar) & \loesung{hablaron} \\
¿tú (viajar)? & \loesung{Viajaste} \\
\end{tabular}

\vspace{1em}

\textbf{Aufgabe 4: Freie Produktion} (8 Punkte)

Bewertungskriterien:
\begin{itemize}
    \item Mindestens 4 Sätze vorhanden: 2P
    \item Signalwort verwendet: 1P
    \item Verschiedene Personen (mind. 2): 2P
    \item Verschiedene Verben (mind. 2): 2P
    \item Sprachliche Korrektheit: 1P
\end{itemize}

% ═══════════════════════════════════════════════════════════════
% 6. HAUSAUFGABE
% ═══════════════════════════════════════════════════════════════

\section*{6. Hausaufgabe}

\begin{itemize}
    \item Lernpfad LP-08c.html bearbeiten (digital)
    \item Vokabeln lernen: Signalwörter + Konjugation hablar
    \item \E{} Optional: Urlaubsbericht erweitern (5--8 Sätze)
\end{itemize}

% ═══════════════════════════════════════════════════════════════
% 7. TAFELBILD
% ═══════════════════════════════════════════════════════════════

\section*{7. Tafelbild (Vorschlag)}

\begin{center}
\fbox{\parbox{0.95\textwidth}{
\textbf{\large El Pretérito Indefinido (-ar)}

\vspace{1em}

\begin{minipage}{0.45\textwidth}
\textbf{Konjugation: hablar}

\begin{tabular}{ll}
yo & habl\textbf{é} \\
tú & habl\textbf{aste} \\
él/ella & habl\textbf{ó} \\
nosotros & habl\textbf{amos} \\
vosotros & habl\textbf{asteis} \\
ellos & habl\textbf{aron} \\
\end{tabular}
\end{minipage}
\hfill
\begin{minipage}{0.45\textwidth}
\textbf{Signalwörter (palabras clave)}

\vspace{0.5em}

\es{ayer} -- gestern

\es{el año pasado} -- letztes Jahr

\es{el verano pasado} -- letzten Sommer

\es{la semana pasada} -- letzte Woche
\end{minipage}

\vspace{1em}

\textbf{Beispiel:} \es{El verano pasado viajé a España.}
}}
\end{center}

% ═══════════════════════════════════════════════════════════════
% 8. REFLEXIONSNOTIZEN
% ═══════════════════════════════════════════════════════════════

\section*{8. Reflexionsnotizen (nach Durchführung)}

\begin{tabularx}{\textwidth}{|l|X|}
\hline
\textbf{Was lief gut?} & \\[2em]
\hline
\textbf{Was lief nicht?} & \\[2em]
\hline
\textbf{Zeitmanagement} & \\[2em]
\hline
\textbf{Für nächstes Mal} & \\[2em]
\hline
\end{tabularx}

\end{document}
